微分方程与动力系统描述系统怎样演化：给定规则与初值，整条轨迹随之确定。控制把这个问题加上一层选择权：方程里留有一个可以由我们实时调节的输入，问题从``系统会去哪''变成``怎样让系统按我们的意愿演化''。意愿必须先写清楚——是``别偏离目标太远''，还是``用最少的时间或燃料到达''——不同的写法通向不同的数学。

最优控制回答模型已知时的情形：状态方程与代价函数都写在纸上，问题被精确地表述为``在时间轴上选一条输入函数，使目标泛函最小''。这时有两条路可走：一是研究最优解必须满足的条件，二是直接构造给出最优策略的函数。本 Part 单元一沿这两条路走，终点处会发现它们在计算上撞上同一堵墙：问题定义在整个状态空间上，维数稍高就解不动。

强化学习面对更窘迫也更现实的处境：连模型都不知道。能用的只有交互——选一个动作，看到下一个状态和一个奖励数字。怎样从这样的试错中学会控制？单元二会看到，答案并不是另起炉灶：把决策问题写成 Markov 决策过程后，学习目标仍然是单元一那个 Bellman 不动点，只是方程里的转移概率与奖励从``已知''变成``只能采样''，解方程随之变成从数据中逼近。

两个单元合起来，主线只有一句话：模型已知时解方程，模型未知时从数据逼近同一个不动点。沿途需要的基础前面已经备好：反馈为什么能稳住系统，用 ``稳定性与 Lyapunov 函数'' 的能量思想；MDP 去掉动作就退回 ``Markov 性质、转移矩阵与多步演化'' 的链；最大值原理导出的两点边值问题，计算上交给 ``边值问题：从打靶到全局离散'' 的方法。
